% Generated draft: requires mathematical and source review.
\documentclass[11pt]{article}
\usepackage[margin=25mm]{geometry}
\usepackage{fontspec}
\setmainfont{Latin Modern Roman}
\usepackage{amsmath,amssymb,mathtools,unicode-math}
\setmathfont{Latin Modern Math}
\usepackage{enumitem,needspace,xurl,hyperref}
\setlength{\parindent}{0pt}
\setlength{\parskip}{5pt}
\setlength{\emergencystretch}{3em}
\newcommand{\N}{\mathbb{N}}
\newcommand{\Z}{\mathbb{Z}}
\newcommand{\Q}{\mathbb{Q}}
\newcommand{\R}{\mathbb{R}}
\newcommand{\C}{\mathbb{C}}
\newcommand{\F}{\mathbb{F}}
\newcommand{\A}{\mathbb{A}}
\newcommand{\PP}{\mathbb{P}}
\newcommand{\E}{\mathbb{E}}
\newcommand{\eps}{\varepsilon}
\newcommand{\dd}{\,\mathrm{d}}
\begin{document}
\section*{UnsolvedMath ID: 10400099}
\textbf{Source problem code:} catalogue entry
\subsection{Definitions and mathematical statement}
For the invariant $h_X$ defined in the omitted preceding source, the conjecture is
\[
\log h_X\text{ is a finite-type (Vassiliev) invariant}
\quad\Longrightarrow\quad \log h_X\text{ is constant}.
\]
Equivalently a nonconstant $\log h_X$ should not have any finite Vassiliev degree. A degree-$\leq d$ invariant extends to singular knots by the crossing-difference rule and vanishes on every knot with $d+1$ transverse double points. The source does not define $h_X$, its domain/normalization, or the hypotheses making its logarithm well defined. It cannot simply be replaced by an arbitrary quandle coloring count, and zero/complex values would require additional logarithm conventions.

The following section 5.3 is separate quandle-cohomology exposition. A quandle has an operation $*$ with $x*x=x$, bijective right translations and $(x*y)*z=(x*z)*(y*z)$. The extract takes $C^r(X;A)$ to be all maps $X^r\to A$ for an additive abelian group $A$, and uses
\[
C_Q^1=C^1,\qquad C_Q^2=\{g:g(x,x)=0\},\qquad
C_Q^3=\{h:h(x,x,y)=h(x,y,y)=0\}.
\]
Its coboundaries are
\[
(d^1f)(x,y)=f(x)-f(x*y),
\]
\[
(d^2g)(x,y,z)=g(x,z)-g(x,y)-g(x*y,z)+g(x*z,y*z),
\]
and $H_Q^2(X;A)=\ker d^2/\operatorname{im}d^1$. The source reports $H_Q^2(X;A)\cong\operatorname{Hom}(H_2^Q(X),A)$ by universal coefficients since $H_1^Q(X)$ is free abelian, with reference [82] for the chain/cochain conventions and general degrees.

The appended Table 5 concerns $H_2^Q,H_3^Q$ of connected quandles of orders three through nine: dihedral examples $R_3,R_5,R_7,R_9$ and the listed Alexander quotients over moduli $2,3,5,7,9$. It also reports $H_2^Q(\Z[t^{\pm1}]/(p,t-a))=0$ for every prime $p$ and $a\ne0,1$ in $\Z/p\Z$. The source attributes the finite table to [264], omits one of each dual pair and nonconnected quandles, and attributes the last family to [284]. The extracted table's row/column alignment and undefined $R_j$/Alexander operation conventions need checking before its remaining group entries can be treated as standalone claims. This contextual table is not a new hypothesis or a substitute for the missing definition of $h_X$.
\subsection{Short English statement}
For the source-defined $h_X$, prove that its logarithm can be a finite-type knot invariant only when constant. The preceding definition is missing; the subsequent quandle-cohomology formulas and table are separate context.
\subsection{Sources}
\begin{itemize}
\item Local UnsolvedMath registry, permanent ID 10400099, exact source code catalogue entry, Conjecture 5.3; no proposer/year supplied. Registry status open, classification OPEN-TRIAGE.
\item Ohtsuki, \emph{Problems on invariants of knots and 3-manifolds} (2002), Conjecture 5.3 and appended section 5.3/Table 5, PDF page 87, \url{https://msp.org/gtm/2002/04/gtm-2002-04-024s.pdf}. Supplied pdf-font-structure extraction accessed July 29, 2026; provenance distinguishes 231 bold numbered items. No PDF was accessed here.
\item Source references [82], [264] and [284] concern quandle (co)homology and table computations; their bibliography entries are absent. The table's omitted dual/nonconnected cases refer to remarks on Problem 5.6.
\item The supplied summary only says the $h_X$ conjecture is not obviously resolved, with no precise later result. Provenance flags current primary-literature status and redistribution terms NEEDS\_REVIEW.
\end{itemize}
\end{document}
